\documentclass[preprint,12pt]{elsarticle}
\usepackage{amsmath,amssymb,amsfonts,bm,mathtools}
\usepackage[a4paper,margin=1in]{geometry}
\usepackage{graphicx,booktabs,array,tabularx,xltabular,multirow,ragged2e,siunitx}
\usepackage[table]{xcolor}
\usepackage[hidelinks]{hyperref}
\newcolumntype{Y}{>{\RaggedRight\arraybackslash}X}
\newcolumntype{L}[1]{>{\RaggedRight\arraybackslash}p{#1}}
\newcommand{\bmk}{\bm{k}}
\newcommand{\barw}{\bar{\omega}}
\journal{International Journal of Mechanical Sciences}

\begin{document}
\begin{frontmatter}
\title{Supplementary Material for: Orientation-Dependent Dispersion, Directional Band Gaps, and Wave Steering in Two-Dimensional Strain-Gradient Elasticity}
\end{frontmatter}

This file contains the complete parameter summary and expanded verification details supporting the main manuscript. External comparisons are distinguished from internal consistency checks; calculations for which the published source data are insufficient are not described as validated benchmarks.

\section*{S1. Parameter summary}
Table~\ref{tab:s1_parameters} lists the material, geometric, and numerical values used in the principal studies and in the external bilayer calculations.
\begin{xltabular}{\textwidth}{@{}l p{3.2cm} Y@{}}
\caption{Complete parameter summary for the homogeneous, composite, and transfer-matrix calculations.}
\label{tab:s1_parameters}\\
\toprule
Study & Parameter group & Values used \\
\midrule
\endfirsthead
\multicolumn{3}{@{}l@{}}{\footnotesize\emph{Table~\ref{tab:s1_parameters} --- continued}}\\
\toprule
Study & Parameter group & Values used \\
\midrule
\endhead
\bottomrule
\endfoot
\bottomrule
\endlastfoot
Case H & Cell and base material & $L=1.0\,\mathrm{m}$; $\lambda=\mu=1.0\,\mathrm{Pa}$; $\rho=1.0\,\mathrm{kg/m^3}$ \\
Case H & Gradient and inertia lengths & $l_{\mathrm{iso}}=0.20\,\mathrm{m}$; $\ell_{\mathrm{i}}=0.20\,\mathrm{m}$ ($\ell_{\mathrm{i}}^2=0.04\,\mathrm{m^2}$) \\
Case H & Anisotropy scaling & Volume-equivalent axes $l_1=l_{\mathrm{iso}}\sqrt{\mathrm{AR}}$ and $l_2=l_{\mathrm{iso}}/\sqrt{\mathrm{AR}}$; $\mathrm{AR}\in\{1,2,3,5,7,10\}$ \\
Case H & Design-map sampling & $\mathrm{AR}\in\{1,2,3,5,7,10\}$ and $\theta\in\{0^\circ,15^\circ,30^\circ,45^\circ,60^\circ,75^\circ,90^\circ\}$; 42 sampled $(\theta,\mathrm{AR})$ points \\
Case H & Frequency-sweep sampling & $\mathrm{AR}=5$: $\theta=0^\circ$--$90^\circ$ in $2.5^\circ$ steps; $\theta=45^\circ$: $\mathrm{AR}\in\{1,2,3,5,7,10\}$ \\
Case H & Steering sampling & $\bar{k}=0.5$; $\mathrm{AR}\in\{5,10\}$; $\theta=0^\circ,15^\circ,\ldots,90^\circ$; ring direction sampled at $5^\circ$ \\
Case H & Band paths and branches & $\Gamma$--$X$--$M$--$\Gamma$, 40 intervals per leg (121 path nodes); four lowest physical branches \\
Case H & Production resolution & $16\times16$ BFS elements per cell ($h/L=1/16$); homogeneous results use the listed $16\times16$ mesh unless stated otherwise \\
Case H & Zone scans & $41\times81$ half-zone grid for complete-gap/IFC evaluations; the Case-C quoted gap uses the separate grid stated below \\
Case H & Mesh-convergence point & $\bmk=(0.31\pi/L,0.22\pi/L)$; $N\times N$, $N\in\{4,8,16,32\}$ \\
Case H & Numerical check tolerances & Hermiticity $<10^{-12}$; symmetry discrepancies $<10^{-10}$; acoustic-slope discrepancy $<10^{-6}$; energy-flux residual $<10^{-6}$ \\
Case C & Epoxy matrix & $\rho_m=1142\,\mathrm{kg/m^3}$; $\mu_m=1.48\,\mathrm{GPa}$; $\lambda_m=4.57\,\mathrm{GPa}$ \cite{zhanwei2010} \\
Case C & YBCO inclusion & $\rho_i=6333\,\mathrm{kg/m^3}$; $\mu_i=37.0\,\mathrm{GPa}$; $\lambda_i=95.0\,\mathrm{GPa}$ \cite{zhanwei2010} \\
Case C & Geometry and gradient lengths & $r_0/L=0.30$ (area fraction $28.27\%$); $\ell_m/L=0.10$; $\ell_i/L=0.20$ \\
Case C & Frequency scale and grid & $c_{t,m}=1138.4\,\mathrm{m/s}$; $\omega_{0,C}=3576.4\,\mathrm{rad/s}$; quoted gap on $4\times4$ mesh and $21\times21$ Brillouin-zone grid \\
B1 & AlN layer A & $a_A=0.01\,\mathrm{m}$; $\rho_A=3230\,\mathrm{kg/m^3}$; $c_{33,A}=3.9\times10^{11}\,\mathrm{Pa}$; $a_{3,A}=8.4\times10^{-11}$ in source units \cite{li2024anchorB} \\
B1 & BaTiO$_3$ layer B & $a_B=0.01\,\mathrm{m}$; $\rho_B=5800\,\mathrm{kg/m^3}$; $c'_{33,B}=1.62\times10^{11}\,\mathrm{Pa}$; $a'_{3,B}=1.26\times10^{-8}$ in source units \cite{li2024anchorB} \\
B1 & Classical setting and normalization & $b=a_A+a_B=0.02\,\mathrm{m}$; $l=l_1=f=0$; $\omega_0=2.2422\times10^6\,\mathrm{rad/s}$ \\
B2 & Published gradient parameters & Printed values $l=10^{-5}$, $l_1=2\times10^{-5}$, $f=F=0$, and layer ratios $L=L_1=5$; the units and dimensional versus normalized interpretation of the lengths are unspecified in the source \cite{li2024anchorB} \\
B2 & Interpretations evaluated & Dimensional micro-scale ($a=l=10^{-5}\,\mathrm{m}$); dimensional 1-cm layer ($a=0.01\,\mathrm{m}$, $l=10^{-5}\,\mathrm{m}$); and normalized-length reading ($\bar l=10^{-5}$). These yield different curves; B2 is not externally validated \\
B3 & Pb/brass material and thickness & $a_1=a_2=10\,\mu\mathrm{m}$; $\mu_1=2.3\times10^{10}\,\mathrm{Pa}$; $\rho_1=7500\,\mathrm{kg/m^3}$; $\mu_2=1.288\times10^9\,\mathrm{Pa}$; $\rho_2=1177.5\,\mathrm{kg/m^3}$ \cite{li2023anchorA} \\
B3 & Gradient, inertia, and frequency scale & $c_1/a_1^2=0.15$; $d_1/a_1=0.25$; $c_2/c_1=d_2/d_1=1.5$; $\omega_{0,B3}=4.114\times10^8\,\mathrm{rad/s}$; the first-layer ratios are inherited from the paper's comparison example, not reported on the plotted panel \\
\end{xltabular}


\section*{S2. External transfer-matrix comparisons}
\subsection*{S2.1. Classical AlN/BaTiO$_3$ bilayer}
The classical bilayer uses $a_A=a_B=0.01\,\mathrm{m}$, so the period is $b=0.02\,\mathrm{m}$. The transfer-matrix trace reduces to the Rytov relation reported in the main text \cite{li2024anchorB,zhengwei2009}. The homogeneous single-layer identity and identical-material bilayer reduction have maximum relative residuals $6.47\times10^{-16}$ and $7.22\times10^{-16}$; the heterogeneous transfer-matrix trace agrees with the Rytov relation to $2.22\times10^{-16}$. The calculated normalized stop-band intervals are $[0.4806,0.5209]$, $[0.9812,1.0215]$, and $[1.4993,1.5028]$. The source comparison is graphical because numerical source curves were not released; no percentage error is inferred from raster plots.

\subsection*{S2.2. Gradient-elastic bilayers and validation limits}
For the AlN/BaTiO$_3$ gradient case, the published plot lists $l=10^{-5}$ and $l_1=2\times10^{-5}$ while the surrounding text gives a macroscale layer width and the plot does not specify whether the length symbols are dimensional or normalized. The paper defines $\bar l=l/b$ with $b=a_A+a_B=0.02\,\mathrm{m}$; reading the printed $10^{-5}$ as $\bar l$ therefore gives $l=2\times10^{-7}\,\mathrm{m}$, whereas using a single layer width as the scale gives $10^{-7}\,\mathrm{m}$. Together with the dimensional-micro and dimensional-macroscale readings, these alternatives yield different transfer-matrix curves. The calculated stop bands for the dimensional-micro interpretation are $[0.078,0.242]$, $[0.378,0.383]$, $[0.742,0.760]$, $[0.785,1.103]$, and $[1.123,1.435]$. The homogeneous-limit and identical-material reductions each have algebraic residual $4.10\times10^{-57}$. These are internal identities, not evidence of agreement with the publication. This case is \emph{not externally validated} because the source convention is unresolved and numerical source data are unavailable \cite{li2024anchorB}.

For the Pb/brass shear-horizontal case, the evaluated ratios are $c_1/a_1^2=0.15$, $d_1/a_1=0.25$, and $c_2/c_1=d_2/d_1=1.5$. The values $c_1$ and $d_1$ are inherited from a shared numerical example in Li et al. rather than specified on the cited comparison panel. The calculated stop-band intervals are $[0.340,1.024]$, $[1.423,1.867]$, and $[2.477,2.911]$. The homogeneous-limit identity and identical-material reduction have relative residuals $1.37\times10^{-13}$ and $4.19\times10^{-14}$, respectively. The implementation agrees with the source's printed dipolar layer-matrix formulation to machine precision; this establishes formulation consistency only. The band curves remain \emph{not externally validated} because the cited panel does not supply its parameter set, normalization, or numerical curve data \cite{li2023anchorA,liweizhou2016}.

\subsection*{S2.3. Pure-bending micro-beam limit}
The micro-beam pure-bending reduction reproduces the closed-form relations reported by Papargyri-Beskou and Beskos to below $10^{-15}$ in both evaluated limits \cite{papargyribeskou2009}. This is an analytical benchmark, separate from the plot-based bilayer comparisons.

\section*{S3. Internal checks and mesh-convergence details}
The main text reports the eight consistency checks and the observed Case-H mesh rate. The expanded diagnostics complement the eight summary checks reported in the main text.

The normalized directional-width values corresponding to the representative Case-H designs in the main-text gap summary are listed in Table~\ref{tab:s3_gap_ratio}. At the reported precision, the selected rows agree between the $8\times8$ and $16\times16$ meshes; over the full 126-point design map the maximum mesh-to-mesh changes are $5.5\times10^{-3}$ for $\Delta_{GX}$, $5.8\times10^{-3}$ for $\Delta_{\mathrm{path}}$, and $5.7\times10^{-3}$ for $\Delta_{\mathrm{complete}}$.
\begin{table}[htbp]
\centering
\caption{Normalized directional-gap widths for the representative Case-H designs.}
\label{tab:s3_gap_ratio}
\begin{tabularx}{\textwidth}{@{}c c r@{}}
\toprule
$\mathrm{AR}$ & $\theta$ & $\Delta/\bar{\omega}_{\mathrm{mid}}$ \\
\midrule
$1$ & $0^\circ$  & $-0.09$ \\
$1$ & $45^\circ$ & $-0.09$ \\
$1$ & $90^\circ$ & $-0.09$ \\
$3$ & $0^\circ$  & $-0.13$ \\
$3$ & $45^\circ$ & $-0.13$ \\
$3$ & $90^\circ$ & $-0.13$ \\
$5$ & $0^\circ$  & $-0.16$ \\
$5$ & $45^\circ$ & $-0.16$ \\
$5$ & $90^\circ$ & $-0.16$ \\
$10$ & $0^\circ$  & $-0.23$ \\
$10$ & $45^\circ$ & $-0.15$ \\
$10$ & $90^\circ$ & $-0.23$ \\
\bottomrule
\end{tabularx}

\end{table}

\begin{itemize}
  \item \textbf{Hermiticity and periodicity:} the stiffness and mass Hermiticity residuals are $2.01\times10^{-16}$ and $5.63\times10^{-17}$. Reciprocal-lattice shifts change the reduced matrices by at most $4.0\times10^{-16}$ and the spectrum by $9.66\times10^{-16}$.
  \item \textbf{Isotropic rotation and rotation--swap symmetry:} at $\mathrm{AR}=1$, the maximum frequency change over the sampled orientations is $9.63\times10^{-16}$. For the anisotropic rotation--swap comparison, the spectral discrepancy is $6.44\times10^{-16}$; omitting the axis swap gives a discrepancy of $0.179$.
  \item \textbf{Acoustic slopes:} the best long-wave speed discrepancy is $1.25\times10^{-8}$. At $\bar{k}=10^{-2}$ the $1.80\times10^{-6}$ deviation includes the expected $O(\bar{k}^2)$ dispersion and is not treated as a discretization error.
  \item \textbf{Definiteness:} the stiffness operator has the two expected translation modes at $\Gamma$; the minimum interior eigenvalue is $\omega^2=3.028$. The minimum mass-matrix eigenvalue is $9.9\times10^{-5}>0$.
  \item \textbf{Micro-inertia and energy flux:} at $\bar{k}=200$, the bounded phase-speed result differs from its asymptote by $2.85\times10^{-4}$. The source manuscript previously recorded a group-velocity/energy-flux residual of $6.91\times10^{-10}$; this check is pending rerun with the corrected harmonic micro-inertial flux term because the required raw harmonic fields are not included.
\end{itemize}

\begin{table}[htbp]
\centering
\caption{Expanded mesh-convergence data, including successive changes and eigensolver reproducibility.}
\label{tab:s3_mesh}
\begin{tabularx}{\textwidth}{@{}c r r r r l@{}}
\toprule
Mesh & $\bar{\omega}_T$ & Relative error & Successive change & Reproducibility & Use in fit \\
\midrule
$4\times4$ & 1.164855406908 & $1.51\times10^{-8}$ & --- & $2.90\times10^{-14}$ & Included \\
$8\times8$ & 1.164855390213 & $7.59\times10^{-10}$ & $1.43\times10^{-8}$ & $1.06\times10^{-13}$ & Included \\
$16\times16$ & 1.164855389383 & $4.63\times10^{-11}$ & $7.12\times10^{-10}$ & $5.92\times10^{-13}$ & Included \\
$32\times32$ & 1.164855389329 & $2.48\times10^{-15}$ & $4.63\times10^{-11}$ & $1.26\times10^{-12}$ & Resolution-limited \\
\bottomrule
\end{tabularx}

\end{table}

The empirical rate $p=4.17$ (95\% CI $[3.15,5.20]$) is fitted to the $4^2$, $8^2$, and $16^2$ levels. The $32^2$ result is below three times its measured reproducibility and is therefore reported as resolution-limited, not used to infer the rate. The operational mesh-change floor is $\varepsilon_\Delta=4.63\times10^{-11}$; it is distinct from the eigensolver reproducibility at the finest level.

\section*{S4. Case-C gap resolution and parameter sensitivity}
Table~\ref{tab:s4_caseC_mesh} separates the finite-element mesh effect from the Brillouin-zone sampling effect for the Case-C complete gap. The four leg/path values at the $4\times4$ mesh are $\Delta[\Gamma\text{--}X]=2.6114$, $\Delta[X\text{--}M]=2.7563$, $\Delta[M\text{--}\Gamma]=3.2871$, and $\Delta[\mathrm{path}]=2.5732$, compared with $\Delta[\mathrm{complete}]=2.5732$. The $X$-point gap overestimates the complete gap on the coarsest mesh because the minimum upper-band frequency lies off $X$ on the $\Gamma$--$M$ diagonal; from $8\times8$ onward the reported $X$ and complete gaps coincide to the shown precision.

\begin{table}[htbp]
\centering
\caption{Case-C complete-gap and $X$-point gap values under mesh refinement.}
\label{tab:s4_caseC_mesh}
\begin{tabularx}{\textwidth}{@{}c r r r r r@{}}
\toprule
Mesh & DOFs & $\Delta_{\mathrm{complete}}$ & $\Delta_X$ & Decrement & Decrement ratio \\
\midrule
$4\times4$ & 200 & 2.5732 & 2.7563 & --- & --- \\
$8\times8$ & 648 & 2.2504 & 2.2504 & 0.3228 & 0.552 \\
$16\times16$ & 2312 & 2.0722 & 2.0722 & 0.17814 & 0.554 \\
$32\times32$ & 8712 & 1.9736 & 1.9736 & 0.09861 & 0.565 \\
$64\times64$ & 33800 & 1.9179 & 1.9179 & 0.05575 & --- \\
\bottomrule
\end{tabularx}

\end{table}

At the $4\times4$ mesh, increasing the Brillouin-zone grid from $11\times11$ to $21\times21$ leaves the complete gap at $2.5732$; the $41\times41$ grid gives $2.5728$. The upper edge changes from $7.1430$ to $7.1426$ as the minimum is better resolved along $\Gamma$--$M$. For $8\times8$ and $16\times16$ meshes, the $11\times11$ and $21\times21$ grids agree to four decimals; for $32\times32$ and $64\times64$ they agree to six decimals. With the FE mesh held at $4\times4$, $4\times4$, $6\times6$, and $8\times8$ Gauss points per cut element give $X$-point gaps of $2.7563$, $2.7226$, and $2.7862$, respectively (changes of $-1.22\%$ and $+1.08\%$ relative to the $4\times4$ rule). Under radius variation, $r_0/L=0.20$--$0.40$ corresponds to filling fractions $12.6\%$--$50.3\%$; the complete gap changes from $0.8288$ ($16.61\%$) at $r_0/L=0.20$ to $2.5732$ ($43.94\%$) at $0.30$, then to $2.4426$ ($37.00\%$) at $0.40$. The mesh-refinement sequence in Table~\ref{tab:s4_caseC_mesh} remains decreasing, so the quoted coarse-mesh maximum is not a mesh-converged continuum value.

\bibliographystyle{elsarticle-num}
\bibliography{references}
\end{document}
